% D5 proof generalization snippet after 049/050

\begin{theorem}[D5 active anticipation chain; checked form, theorem target]
Let $m$ be an odd modulus and consider the active best-seed branch of the
unresolved channel $R1\to H_{L1}$. Let
\[
B=(s,u,v,\lambda,f),
\qquad
c=\mathbf 1_{\{q=m-1\}},
\qquad
d=\mathbf 1_{\{U^+\ge m-3\}}.
\]
On the checked odd moduli $m=5,7,9,11$, the branch factors as
\[
B \leftarrow B+c \leftarrow B+c+d.
\]
Moreover:
\begin{enumerate}
    \item the exceptional trigger descends to $B$;
    \item the regular trigger descends to $B+c$ but not to $B$;
    \item deterministic active evolution closes on $B+c+d$;
    \item the carry sheet $c$ is not a current grouped-state observable, but a
    one-sided anticipation datum over the grouped base.
\end{enumerate}
\end{theorem}

\begin{theorem}[Carry as anticipation datum; conceptual theorem]
Let $\tau$ denote the initial flat-run length of the future grouped-delta
sequence, and let $\epsilon_4\in\{\mathrm{flat},\mathrm{wrap},\mathrm{carry\_jump},\mathrm{other}\}$
be the current grouped-delta event class. Then on the checked active branch for $m=5,7,9,11$,
the carry sheet $c$ is a function of $(B,\tau,\epsilon_4)$.

More precisely, all ambiguity of $(B,\tau)$ lies at the boundary $\tau=0$,
and the boundary class is genuinely three-way:
\[
\{\mathrm{wrap},\mathrm{carry\_jump},\mathrm{other}\}.
\]
Thus the canonical theorem-side object is $(B,\tau,\epsilon_4)$.
\end{theorem}

\begin{proposition}[Countdown carrier reduction]
Define $\tau(x)$ to be the length of the initial flat prefix of the future
$dn$-sequence at $x$. Then
\[
\tau(Fx)=\tau(x)-1 \qquad \text{whenever } \tau(x)>0.
\]
Hence all nontrivial dynamics of $\tau$ are confined to the boundary $\tau=0$.
On the checked odd moduli through $m=19$, the boundary reset law is exact on
small current quotients:
\[
R(\epsilon_4,s,u,v,\lambda)=
\begin{cases}
0, & \epsilon_4=\mathrm{wrap},\\
R_{\mathrm{cj}}(s,v,\lambda), & \epsilon_4=\mathrm{carry\_jump},\\
R_{\mathrm{oth}}(s,u,\lambda), & \epsilon_4=\mathrm{other},
\end{cases}
\]
with observed value sets
\[
R_{\mathrm{cj}}\subseteq\{0,1,m-2\},
\qquad
R_{\mathrm{oth}}\subseteq\{0,m-4,m-3\}.
\]
Consequently, any realization of the D5 carry mechanism may be reformulated as
realization of a countdown carrier together with a boundary reset map.
\end{proposition}

\begin{proposition}[Witness family for bounded-horizon lower bounds]
For each checked odd modulus $m$, there are states
\[
x^-_m=(m-2,2,1,2,0),
\qquad
x^+_m=(m-1,2,1,2,0),
\]
with common grouped base
\[
B_*=(3,1,2,0,\mathrm{regular}),
\]
common current event class $\epsilon_4=\mathrm{flat}$, but opposite carry
labels and anticipation values
\[
c(x^-_m)=0,
\quad
c(x^+_m)=1,
\qquad
\tau(x^-_m)=m-3,
\quad
\tau(x^+_m)=m-4.
\]
They also share a common future flat/nonflat prefix of length $m-5$.
Therefore any coding that sees only current $B$, current $\epsilon_4$, and the
next $h<m-4$ future flat/nonflat bits fails on modulus $m$.
\end{proposition}

\begin{remark}[Source-residue refinement]
The stronger checked coordinate
\[
\rho=u_{\mathrm{source}}+1 \pmod m
\]
may still be useful as a constructive refinement: on the checked odd moduli
through $m=19$, $\tau$ is exact on $(s,u,v,\lambda,\rho)$, $c$ is exact on
$(u,\rho,\epsilon_4)$, and
\[
q \equiv u-\rho + \mathbf 1_{\{\epsilon_4=\mathrm{carry\_jump}\}} \pmod m.
\]
However, for checked $m\ge 7$, $\rho$ is not recoverable from
$(B,\tau,\epsilon_4)$, so it should be treated as a constructive refinement
rather than as the minimal theorem-side coordinate.
\end{remark}
